\documentclass[11pt]{article}

\usepackage[margin=2.2cm]{geometry}
\usepackage{amsmath,amssymb}
\usepackage{siunitx}
\usepackage{booktabs}
\usepackage{longtable}
\usepackage{array}
\usepackage[table]{xcolor}
\usepackage{caption}
\usepackage{hyperref}
\usepackage{titlesec}

\sisetup{output-exponent-marker=\ensuremath{\mathrm{e}}, retain-explicit-plus=false, group-digits=false}

\titlespacing*{\section}{0pt}{14pt}{6pt}
\titlespacing*{\subsection}{0pt}{10pt}{4pt}

\definecolor{hdrblue}{HTML}{2C3E50}
\definecolor{rowgray}{HTML}{F2F2F2}

\newcommand{\hdr}[1]{\textcolor{white}{\textbf{#1}}}

\title{\textbf{CaF--He Collision Estimate: Formulas and Distance Tables}}
\author{by \emph{Claude}}
\date{\today}

\begin{document}
\maketitle
\vspace{-2.5em}

\noindent
Estimate of the number of CaF--He collisions from the buffer-gas-cell nozzle exit ($z=0$) out to
$z = 1\,\text{m}$, for a cryogenic buffer-gas source adapted for pulsed-ablation CaF production on the
continuous-flow cell design of Gangwar \textit{et al.}, ``Buffer gas cooling of a continuous CO
molecular beam,'' \textit{Optics Express} \textbf{33}, 28164 (2025),
\href{https://doi.org/10.1364/OE.560951}{doi:10.1364/OE.560951}.

\section*{Input parameters}

\begin{center}
\begin{tabular}{@{}p{3.3cm}p{1.4cm}p{3.0cm}p{6.3cm}@{}}
\toprule
\rowcolor{hdrblue}
\hdr{Quantity} & \hdr{Symbol} & \hdr{Value} & \hdr{Source} \\
\midrule
Cell / nozzle temperature & $T$ & \SI{4}{K} & user-specified (cell 4 K stage; 3.7 K measured in the paper) \\
\rowcolor{rowgray} He flow rate & $f_{He}$ & 20 sccm $= \num{9.0e18}\,\text{s}^{-1}$ & user-specified; 1 sccm $=\num{4.5e17}\,\text{s}^{-1}$ (paper's own conversion) \\
Nozzle diameter & $d$ & 3 mm and 5 mm & range tested in the paper (2--5 mm); brackets typical CaF cells \\
\rowcolor{rowgray} CaF--He cross section & $\sigma$ & $\num{1e-14}\,\text{cm}^2$ & order-of-magnitude value used throughout buffer-gas literature (paper uses the same value for CO--He) \\
Background chamber pressure & $P_{bg}$ & $\num{1e-7}$ mbar & paper's stated typical operating pressure, both chambers \\
\rowcolor{rowgray} Background gas temperature & $T_{bg}$ & \SI{300}{K} & gauge reads pressure referenced to room temperature \\
\bottomrule
\end{tabular}
\end{center}

\vspace{0.3cm}
\noindent Two collision zones are combined: (1) a \emph{near-field} zone right at the nozzle where the
He is still dense and the flow is expanding hydrodynamically (this is where essentially all the
collisions happen), and (2) a \emph{far-field} zone where the molecule coasts through the low residual
chamber pressure. Collisions inside the buffer-gas cell itself (thermalization, upstream of $z=0$) are
not included here.

\section*{Formulas used}

\subsection*{1. Mean thermal speed of He at temperature $T$}
\begin{equation}
\bar{v}_{He} = \sqrt{\dfrac{8 k_B T}{\pi m_{He}}}
\end{equation}

\subsection*{2. He number density at the aperture (stagnation density)}
\begin{equation}
n_0 = \dfrac{4 f_{He}}{A_{aperture}\, \bar{v}_{He}}, \qquad
A_{aperture} = \dfrac{\pi d^2}{4}
\end{equation}
where $f_{He}$ is the He flow in particles/s (1 sccm $= \num{4.5e17}\,\text{s}^{-1}$).

\subsection*{3. He--He mean free path near the aperture}
\begin{equation}
\lambda_{He} = \dfrac{1}{\sqrt{2}\, n_0\, \sigma_{He\text{-}He}}
\end{equation}

\subsection*{4. Reynolds number of the expansion}
\begin{equation}
Re \approx \dfrac{2 d}{\lambda_{He}}
\end{equation}

\subsection*{5. Total collisions picked up near the orifice}
\begin{equation}
N_{\text{near-field}} \approx Re/2
\end{equation}
(standard buffer-gas-beam result, Hutzler \textit{et al.}, \textit{Chem.\ Rev.} \textbf{112}, 4803 (2012);
used directly in the source paper.)

\subsection*{6. Local He density vs.\ distance downstream (free-jet scaling)}
\begin{equation}
n(z) =
\begin{cases}
n_0, & z \le z^{*} \\[4pt]
0.15\, n_0 \left(\dfrac{d}{z}\right)^{2}, & z > z^{*}
\end{cases}
\qquad\text{floored at } n_{bg}\ \text{if smaller,} \qquad z^{*} = \sqrt{0.15}\; d
\end{equation}

\subsection*{7. Local collision rate (inverse local mean free path)}
\begin{equation}
\dfrac{dN}{dz}(z) = n(z)\, \sigma_{CaF\text{-}He} \qquad [\text{cm}^{-1}]
\end{equation}

\subsection*{8. Cumulative collisions from nozzle exit to distance $z$}
\begin{equation}
N(z) = \int_0^{z} \dfrac{dN}{dz'}\, dz' \;=\; N_{\text{near-field}} \cdot f(z) \;+\; n_{bg}\, \sigma_{CaF\text{-}He}\, z
\end{equation}
\begin{equation}
f(z) = \dfrac{z}{2 z^{*}} \ (z \le z^{*}), \qquad
f(z) = 1 - \dfrac{z^{*}}{2 z} \ (z > z^{*})
\end{equation}
where $n_{bg} = P_{bg} / (k_B T_{bg})$ is the background chamber density from the gauge reading.

\subsection*{9. Background (far-field) contribution}
\begin{equation}
N_{bg}(z) = n_{bg}\, \sigma_{CaF\text{-}He}\, z
\end{equation}

\bigskip
\noindent\textit{Note: the near-field density profile $n(z)$ (item 6) is a standard free-jet-expansion
scaling ($z^{-2}$ falloff past a few nozzle diameters), not itself taken from the paper; it is normalized
so that its integral reproduces the paper's cited near-orifice total, $N_{\text{near-field}}\approx Re/2$.
The $z^{-2}$ shape and the $Re/2$ total are therefore on different footings: the total is a
literature-standard result, the distance-dependence is a reasonable modeling choice calibrated to it.}

\clearpage
\section*{Derived quantities for the two nozzle sizes}

\begin{center}
\begin{tabular}{@{}lcc@{}}
\toprule
\rowcolor{hdrblue}
\hdr{Quantity} & \hdr{3 mm nozzle} & \hdr{5 mm nozzle} \\
\midrule
Aperture area $A$ (cm$^2$) & \num{0.07069} & \num{0.19635} \\
\rowcolor{rowgray} Cell/aperture density $n_0$ (cm$^{-3}$) & \num{3.501e16} & \num{1.2605e16} \\
He--He mean free path ($\mu$m) & \num{20.20} & \num{56.10} \\
\rowcolor{rowgray} Reynolds number $Re$ & \num{297.1} & \num{178.3} \\
Near-field total, $Re/2$ & \num{148.5} & \num{89.13} \\
\rowcolor{rowgray} Matching distance $z^{*}$ (mm) & \num{1.162} & \num{1.936} \\
Background density $n_{bg}$ (cm$^{-3}$) & \num{2.414e9} & \num{2.414e9} \\
\bottomrule
\end{tabular}
\end{center}

\clearpage
\section*{Results: 3 mm nozzle}

\subsection*{Fine-resolution table, 0--10 mm (0.5 mm steps)}
\noindent\textit{This is where essentially all the near-field collisions occur.}

\begin{longtable}{@{}>{\raggedleft\arraybackslash}p{1.6cm} >{\raggedleft\arraybackslash}p{3.6cm} >{\raggedleft\arraybackslash}p{3.6cm} >{\raggedleft\arraybackslash}p{3cm}@{}}
\toprule
\rowcolor{hdrblue}
\hdr{$z$ (mm)} & \hdr{rate $dN/dz$ (cm$^{-1}$)} & \hdr{collisions in 0.5-mm bin} & \hdr{cumulative $N(z)$} \\
\midrule
\endfirsthead
\toprule
\rowcolor{hdrblue}
\hdr{$z$ (mm)} & \hdr{rate $dN/dz$ (cm$^{-1}$)} & \hdr{collisions in 0.5-mm bin} & \hdr{cumulative $N(z)$} \\
\midrule
\endhead
\bottomrule
\endfoot
0 & \num{350.126} & \num{0} & \num{0} \\
0.5 & \num{350.126} & \num{31.962} & \num{31.962} \\
1 & \num{350.126} & \num{31.962} & \num{63.9239} \\
1.5 & \num{210.076} & \num{27.0903} & \num{91.0143} \\
2 & \num{118.167} & \num{14.3829} & \num{105.397} \\
2.5 & \num{75.6272} & \num{8.62973} & \num{114.027} \\
3 & \num{52.5189} & \num{5.75316} & \num{119.78} \\
3.5 & \num{38.5853} & \num{4.1094} & \num{123.889} \\
4 & \num{29.5419} & \num{3.08205} & \num{126.971} \\
4.5 & \num{23.3417} & \num{2.39715} & \num{129.369} \\
5 & \num{18.9068} & \num{1.91772} & \num{131.286} \\
5.5 & \num{15.6255} & \num{1.56904} & \num{132.855} \\
6 & \num{13.1297} & \num{1.30754} & \num{134.163} \\
6.5 & \num{11.1875} & \num{1.10638} & \num{135.269} \\
7 & \num{9.64632} & \num{0.948323} & \num{136.218} \\
7.5 & \num{8.40302} & \num{0.82188} & \num{137.04} \\
8 & \num{7.38547} & \num{0.719146} & \num{137.759} \\
8.5 & \num{6.54214} & \num{0.63454} & \num{138.393} \\
9 & \num{5.83543} & \num{0.564036} & \num{138.957} \\
9.5 & \num{5.23734} & \num{0.504664} & \num{139.462} \\
10 & \num{4.7267} & \num{0.454198} & \num{139.916} \\
\end{longtable}

\subsection*{Per-centimeter table, 0--100 cm (1 cm steps)}
\noindent\textit{Rate and per-bin collisions collapse to the residual background-gas value beyond a few
cm; cumulative total asymptotes to $Re/2 = 148.5$.}

\begin{longtable}{@{}>{\raggedleft\arraybackslash}p{1.6cm} >{\raggedleft\arraybackslash}p{3.6cm} >{\raggedleft\arraybackslash}p{3.6cm} >{\raggedleft\arraybackslash}p{3cm}@{}}
\toprule
\rowcolor{hdrblue}
\hdr{$z$ (cm)} & \hdr{rate $dN/dz$ (cm$^{-1}$)} & \hdr{collisions in 1-cm bin} & \hdr{cumulative $N(z)$} \\
\midrule
\endfirsthead
\toprule
\rowcolor{hdrblue}
\hdr{$z$ (cm)} & \hdr{rate $dN/dz$ (cm$^{-1}$)} & \hdr{collisions in 1-cm bin} & \hdr{cumulative $N(z)$} \\
\midrule
\endhead
\bottomrule
\endfoot
0 & \num{350.126} & \num{0} & \num{0} \\
1 & \num{4.7267} & \num{139.916} & \num{139.916} \\
2 & \num{1.18167} & \num{4.31489} & \num{144.231} \\
3 & \num{0.525189} & \num{1.43831} & \num{145.669} \\
4 & \num{0.295419} & \num{0.719168} & \num{146.388} \\
5 & \num{0.189068} & \num{0.431511} & \num{146.82} \\
6 & \num{0.131297} & \num{0.287682} & \num{147.108} \\
7 & \num{0.0964632} & \num{0.205494} & \num{147.313} \\
8 & \num{0.0738547} & \num{0.154127} & \num{147.467} \\
9 & \num{0.0583543} & \num{0.119882} & \num{147.587} \\
10 & \num{0.047267} & \num{0.0959101} & \num{147.683} \\
11 & \num{0.0390636} & \num{0.0784763} & \num{147.762} \\
12 & \num{0.0328243} & \num{0.0654009} & \num{147.827} \\
13 & \num{0.0279686} & \num{0.0553429} & \num{147.882} \\
14 & \num{0.0241158} & \num{0.0474403} & \num{147.93} \\
15 & \num{0.0210076} & \num{0.0411181} & \num{147.971} \\
16 & \num{0.0184637} & \num{0.0359814} & \num{148.007} \\
17 & \num{0.0163554} & \num{0.0317511} & \num{148.039} \\
18 & \num{0.0145886} & \num{0.0282259} & \num{148.067} \\
19 & \num{0.0130933} & \num{0.0252573} & \num{148.092} \\
20 & \num{0.0118167} & \num{0.022734} & \num{148.115} \\
21 & \num{0.0107181} & \num{0.0205711} & \num{148.135} \\
22 & \num{0.00976591} & \num{0.0187032} & \num{148.154} \\
23 & \num{0.00893516} & \num{0.0170789} & \num{148.171} \\
24 & \num{0.00820607} & \num{0.0156577} & \num{148.187} \\
25 & \num{0.00756272} & \num{0.014407} & \num{148.201} \\
26 & \num{0.00699216} & \num{0.0133007} & \num{148.215} \\
27 & \num{0.00648381} & \num{0.0123172} & \num{148.227} \\
28 & \num{0.00602895} & \num{0.0114391} & \num{148.238} \\
29 & \num{0.00562033} & \num{0.0106519} & \num{148.249} \\
30 & \num{0.00525189} & \num{0.00994338} & \num{148.259} \\
31 & \num{0.00491852} & \num{0.00930342} & \num{148.268} \\
32 & \num{0.00461592} & \num{0.00872347} & \num{148.277} \\
33 & \num{0.0043404} & \num{0.00819624} & \num{148.285} \\
34 & \num{0.00408884} & \num{0.00771553} & \num{148.293} \\
35 & \num{0.00385853} & \num{0.00727602} & \num{148.3} \\
36 & \num{0.00364714} & \num{0.00687314} & \num{148.307} \\
37 & \num{0.00345267} & \num{0.00650292} & \num{148.313} \\
38 & \num{0.00327334} & \num{0.00616193} & \num{148.32} \\
39 & \num{0.00310763} & \num{0.00584717} & \num{148.325} \\
40 & \num{0.00295419} & \num{0.00555602} & \num{148.331} \\
41 & \num{0.00281184} & \num{0.00528617} & \num{148.336} \\
42 & \num{0.00267953} & \num{0.0050356} & \num{148.341} \\
43 & \num{0.00255635} & \num{0.00480251} & \num{148.346} \\
44 & \num{0.00244148} & \num{0.00458531} & \num{148.351} \\
45 & \num{0.00233417} & \num{0.00438259} & \num{148.355} \\
46 & \num{0.00223379} & \num{0.0041931} & \num{148.359} \\
47 & \num{0.00213975} & \num{0.00401569} & \num{148.363} \\
48 & \num{0.00205152} & \num{0.00384938} & \num{148.367} \\
49 & \num{0.00196864} & \num{0.00369325} & \num{148.371} \\
50 & \num{0.00189068} & \num{0.00354648} & \num{148.374} \\
51 & \num{0.00181726} & \num{0.00340835} & \num{148.378} \\
52 & \num{0.00174804} & \num{0.00327819} & \num{148.381} \\
53 & \num{0.0016827} & \num{0.0031554} & \num{148.384} \\
54 & \num{0.00162095} & \num{0.00303942} & \num{148.387} \\
55 & \num{0.00156255} & \num{0.00292978} & \num{148.39} \\
56 & \num{0.00150724} & \num{0.002826} & \num{148.393} \\
57 & \num{0.00145482} & \num{0.00272769} & \num{148.396} \\
58 & \num{0.00140508} & \num{0.00263447} & \num{148.398} \\
59 & \num{0.00135786} & \num{0.00254598} & \num{148.401} \\
60 & \num{0.00131297} & \num{0.00246192} & \num{148.403} \\
61 & \num{0.00127028} & \num{0.00238199} & \num{148.406} \\
62 & \num{0.00122963} & \num{0.00230593} & \num{148.408} \\
63 & \num{0.0011909} & \num{0.0022335} & \num{148.41} \\
64 & \num{0.00115398} & \num{0.00216445} & \num{148.413} \\
65 & \num{0.00111875} & \num{0.0020986} & \num{148.415} \\
66 & \num{0.0010851} & \num{0.00203574} & \num{148.417} \\
67 & \num{0.00105295} & \num{0.00197569} & \num{148.419} \\
68 & \num{0.00102221} & \num{0.00191829} & \num{148.421} \\
69 & \num{0.000992795} & \num{0.00186339} & \num{148.422} \\
70 & \num{0.000964632} & \num{0.00181084} & \num{148.424} \\
71 & \num{0.000937651} & \num{0.00176051} & \num{148.426} \\
72 & \num{0.000911786} & \num{0.00171228} & \num{148.428} \\
73 & \num{0.000886977} & \num{0.00166603} & \num{148.429} \\
74 & \num{0.000863166} & \num{0.00162165} & \num{148.431} \\
75 & \num{0.000840302} & \num{0.00157905} & \num{148.433} \\
76 & \num{0.000818334} & \num{0.00153813} & \num{148.434} \\
77 & \num{0.000797217} & \num{0.00149881} & \num{148.436} \\
78 & \num{0.000776906} & \num{0.001461} & \num{148.437} \\
79 & \num{0.000757362} & \num{0.00142462} & \num{148.438} \\
80 & \num{0.000738547} & \num{0.00138961} & \num{148.44} \\
81 & \num{0.000720424} & \num{0.00135589} & \num{148.441} \\
82 & \num{0.000702959} & \num{0.00132341} & \num{148.443} \\
83 & \num{0.000686123} & \num{0.0012921} & \num{148.444} \\
84 & \num{0.000669884} & \num{0.00126191} & \num{148.445} \\
85 & \num{0.000654214} & \num{0.00123279} & \num{148.446} \\
86 & \num{0.000639089} & \num{0.00120468} & \num{148.448} \\
87 & \num{0.000624481} & \num{0.00117754} & \num{148.449} \\
88 & \num{0.000610369} & \num{0.00115133} & \num{148.45} \\
89 & \num{0.00059673} & \num{0.001126} & \num{148.451} \\
90 & \num{0.000583543} & \num{0.00110151} & \num{148.452} \\
91 & \num{0.000570788} & \num{0.00107783} & \num{148.453} \\
92 & \num{0.000558447} & \num{0.00105493} & \num{148.454} \\
93 & \num{0.000546502} & \num{0.00103276} & \num{148.455} \\
94 & \num{0.000534936} & \num{0.0010113} & \num{148.456} \\
95 & \num{0.000523734} & \num{0.000990519} & \num{148.457} \\
96 & \num{0.00051288} & \num{0.000970386} & \num{148.458} \\
97 & \num{0.000502359} & \num{0.000950876} & \num{148.459} \\
98 & \num{0.000492159} & \num{0.000931963} & \num{148.46} \\
99 & \num{0.000482267} & \num{0.000913623} & \num{148.461} \\
100 & \num{0.00047267} & \num{0.000895833} & \num{148.462} \\
\end{longtable}

\clearpage
\section*{Results: 5 mm nozzle}

\subsection*{Fine-resolution table, 0--10 mm (0.5 mm steps)}
\noindent\textit{This is where essentially all the near-field collisions occur.}

\begin{longtable}{@{}>{\raggedleft\arraybackslash}p{1.6cm} >{\raggedleft\arraybackslash}p{3.6cm} >{\raggedleft\arraybackslash}p{3.6cm} >{\raggedleft\arraybackslash}p{3cm}@{}}
\toprule
\rowcolor{hdrblue}
\hdr{$z$ (mm)} & \hdr{rate $dN/dz$ (cm$^{-1}$)} & \hdr{collisions in 0.5-mm bin} & \hdr{cumulative $N(z)$} \\
\midrule
\endfirsthead
\toprule
\rowcolor{hdrblue}
\hdr{$z$ (mm)} & \hdr{rate $dN/dz$ (cm$^{-1}$)} & \hdr{collisions in 0.5-mm bin} & \hdr{cumulative $N(z)$} \\
\midrule
\endhead
\bottomrule
\endfoot
0 & \num{126.045} & \num{0} & \num{0} \\
0.5 & \num{126.045} & \num{11.5063} & \num{11.5063} \\
1 & \num{126.045} & \num{11.5063} & \num{23.0126} \\
1.5 & \num{126.045} & \num{11.5063} & \num{34.5189} \\
2 & \num{118.167} & \num{11.4599} & \num{45.9788} \\
2.5 & \num{75.6272} & \num{8.62973} & \num{54.6086} \\
3 & \num{52.5189} & \num{5.75316} & \num{60.3617} \\
3.5 & \num{38.5853} & \num{4.1094} & \num{64.4711} \\
4 & \num{29.5419} & \num{3.08205} & \num{67.5532} \\
4.5 & \num{23.3417} & \num{2.39715} & \num{69.9503} \\
5 & \num{18.9068} & \num{1.91772} & \num{71.868} \\
5.5 & \num{15.6255} & \num{1.56904} & \num{73.4371} \\
6 & \num{13.1297} & \num{1.30754} & \num{74.7446} \\
6.5 & \num{11.1875} & \num{1.10638} & \num{75.851} \\
7 & \num{9.64632} & \num{0.948323} & \num{76.7993} \\
7.5 & \num{8.40302} & \num{0.82188} & \num{77.6212} \\
8 & \num{7.38547} & \num{0.719146} & \num{78.3403} \\
8.5 & \num{6.54214} & \num{0.63454} & \num{78.9749} \\
9 & \num{5.83543} & \num{0.564036} & \num{79.5389} \\
9.5 & \num{5.23734} & \num{0.504664} & \num{80.0436} \\
10 & \num{4.7267} & \num{0.454198} & \num{80.4978} \\
\end{longtable}

\subsection*{Per-centimeter table, 0--100 cm (1 cm steps)}
\noindent\textit{Rate and per-bin collisions collapse to the residual background-gas value beyond a few
cm; cumulative total asymptotes to $Re/2 = 89.13$.}

\begin{longtable}{@{}>{\raggedleft\arraybackslash}p{1.6cm} >{\raggedleft\arraybackslash}p{3.6cm} >{\raggedleft\arraybackslash}p{3.6cm} >{\raggedleft\arraybackslash}p{3cm}@{}}
\toprule
\rowcolor{hdrblue}
\hdr{$z$ (cm)} & \hdr{rate $dN/dz$ (cm$^{-1}$)} & \hdr{collisions in 1-cm bin} & \hdr{cumulative $N(z)$} \\
\midrule
\endfirsthead
\toprule
\rowcolor{hdrblue}
\hdr{$z$ (cm)} & \hdr{rate $dN/dz$ (cm$^{-1}$)} & \hdr{collisions in 1-cm bin} & \hdr{cumulative $N(z)$} \\
\midrule
\endhead
\bottomrule
\endfoot
0 & \num{126.045} & \num{0} & \num{0} \\
1 & \num{4.7267} & \num{80.4978} & \num{80.4978} \\
2 & \num{1.18167} & \num{4.31489} & \num{84.8127} \\
3 & \num{0.525189} & \num{1.43831} & \num{86.251} \\
4 & \num{0.295419} & \num{0.719168} & \num{86.9702} \\
5 & \num{0.189068} & \num{0.431511} & \num{87.4017} \\
6 & \num{0.131297} & \num{0.287682} & \num{87.6893} \\
7 & \num{0.0964632} & \num{0.205494} & \num{87.8948} \\
8 & \num{0.0738547} & \num{0.154127} & \num{88.049} \\
9 & \num{0.0583543} & \num{0.119882} & \num{88.1688} \\
10 & \num{0.047267} & \num{0.0959101} & \num{88.2648} \\
11 & \num{0.0390636} & \num{0.0784763} & \num{88.3432} \\
12 & \num{0.0328243} & \num{0.0654009} & \num{88.4086} \\
13 & \num{0.0279686} & \num{0.0553429} & \num{88.464} \\
14 & \num{0.0241158} & \num{0.0474403} & \num{88.5114} \\
15 & \num{0.0210076} & \num{0.0411181} & \num{88.5525} \\
16 & \num{0.0184637} & \num{0.0359814} & \num{88.5885} \\
17 & \num{0.0163554} & \num{0.0317511} & \num{88.6203} \\
18 & \num{0.0145886} & \num{0.0282259} & \num{88.6485} \\
19 & \num{0.0130933} & \num{0.0252573} & \num{88.6738} \\
20 & \num{0.0118167} & \num{0.022734} & \num{88.6965} \\
21 & \num{0.0107181} & \num{0.0205711} & \num{88.7171} \\
22 & \num{0.00976591} & \num{0.0187032} & \num{88.7358} \\
23 & \num{0.00893516} & \num{0.0170789} & \num{88.7528} \\
24 & \num{0.00820607} & \num{0.0156577} & \num{88.7685} \\
25 & \num{0.00756272} & \num{0.014407} & \num{88.7829} \\
26 & \num{0.00699216} & \num{0.0133007} & \num{88.7962} \\
27 & \num{0.00648381} & \num{0.0123172} & \num{88.8085} \\
28 & \num{0.00602895} & \num{0.0114391} & \num{88.82} \\
29 & \num{0.00562033} & \num{0.0106519} & \num{88.8306} \\
30 & \num{0.00525189} & \num{0.00994338} & \num{88.8406} \\
31 & \num{0.00491852} & \num{0.00930342} & \num{88.8499} \\
32 & \num{0.00461592} & \num{0.00872347} & \num{88.8586} \\
33 & \num{0.0043404} & \num{0.00819624} & \num{88.8668} \\
34 & \num{0.00408884} & \num{0.00771553} & \num{88.8745} \\
35 & \num{0.00385853} & \num{0.00727602} & \num{88.8818} \\
36 & \num{0.00364714} & \num{0.00687314} & \num{88.8886} \\
37 & \num{0.00345267} & \num{0.00650292} & \num{88.8951} \\
38 & \num{0.00327334} & \num{0.00616193} & \num{88.9013} \\
39 & \num{0.00310763} & \num{0.00584717} & \num{88.9072} \\
40 & \num{0.00295419} & \num{0.00555602} & \num{88.9127} \\
41 & \num{0.00281184} & \num{0.00528617} & \num{88.918} \\
42 & \num{0.00267953} & \num{0.0050356} & \num{88.923} \\
43 & \num{0.00255635} & \num{0.00480251} & \num{88.9278} \\
44 & \num{0.00244148} & \num{0.00458531} & \num{88.9324} \\
45 & \num{0.00233417} & \num{0.00438259} & \num{88.9368} \\
46 & \num{0.00223379} & \num{0.0041931} & \num{88.941} \\
47 & \num{0.00213975} & \num{0.00401569} & \num{88.945} \\
48 & \num{0.00205152} & \num{0.00384938} & \num{88.9489} \\
49 & \num{0.00196864} & \num{0.00369325} & \num{88.9526} \\
50 & \num{0.00189068} & \num{0.00354648} & \num{88.9561} \\
51 & \num{0.00181726} & \num{0.00340835} & \num{88.9595} \\
52 & \num{0.00174804} & \num{0.00327819} & \num{88.9628} \\
53 & \num{0.0016827} & \num{0.0031554} & \num{88.9659} \\
54 & \num{0.00162095} & \num{0.00303942} & \num{88.969} \\
55 & \num{0.00156255} & \num{0.00292978} & \num{88.9719} \\
56 & \num{0.00150724} & \num{0.002826} & \num{88.9747} \\
57 & \num{0.00145482} & \num{0.00272769} & \num{88.9775} \\
58 & \num{0.00140508} & \num{0.00263447} & \num{88.9801} \\
59 & \num{0.00135786} & \num{0.00254598} & \num{88.9826} \\
60 & \num{0.00131297} & \num{0.00246192} & \num{88.9851} \\
61 & \num{0.00127028} & \num{0.00238199} & \num{88.9875} \\
62 & \num{0.00122963} & \num{0.00230593} & \num{88.9898} \\
63 & \num{0.0011909} & \num{0.0022335} & \num{88.992} \\
64 & \num{0.00115398} & \num{0.00216445} & \num{88.9942} \\
65 & \num{0.00111875} & \num{0.0020986} & \num{88.9963} \\
66 & \num{0.0010851} & \num{0.00203574} & \num{88.9983} \\
67 & \num{0.00105295} & \num{0.00197569} & \num{89.0003} \\
68 & \num{0.00102221} & \num{0.00191829} & \num{89.0022} \\
69 & \num{0.000992795} & \num{0.00186339} & \num{89.0041} \\
70 & \num{0.000964632} & \num{0.00181084} & \num{89.0059} \\
71 & \num{0.000937651} & \num{0.00176051} & \num{89.0077} \\
72 & \num{0.000911786} & \num{0.00171228} & \num{89.0094} \\
73 & \num{0.000886977} & \num{0.00166603} & \num{89.011} \\
74 & \num{0.000863166} & \num{0.00162165} & \num{89.0127} \\
75 & \num{0.000840302} & \num{0.00157905} & \num{89.0142} \\
76 & \num{0.000818334} & \num{0.00153813} & \num{89.0158} \\
77 & \num{0.000797217} & \num{0.00149881} & \num{89.0173} \\
78 & \num{0.000776906} & \num{0.001461} & \num{89.0187} \\
79 & \num{0.000757362} & \num{0.00142462} & \num{89.0202} \\
80 & \num{0.000738547} & \num{0.00138961} & \num{89.0215} \\
81 & \num{0.000720424} & \num{0.00135589} & \num{89.0229} \\
82 & \num{0.000702959} & \num{0.00132341} & \num{89.0242} \\
83 & \num{0.000686123} & \num{0.0012921} & \num{89.0255} \\
84 & \num{0.000669884} & \num{0.00126191} & \num{89.0268} \\
85 & \num{0.000654214} & \num{0.00123279} & \num{89.028} \\
86 & \num{0.000639089} & \num{0.00120468} & \num{89.0292} \\
87 & \num{0.000624481} & \num{0.00117754} & \num{89.0304} \\
88 & \num{0.000610369} & \num{0.00115133} & \num{89.0315} \\
89 & \num{0.00059673} & \num{0.001126} & \num{89.0327} \\
90 & \num{0.000583543} & \num{0.00110151} & \num{89.0338} \\
91 & \num{0.000570788} & \num{0.00107783} & \num{89.0349} \\
92 & \num{0.000558447} & \num{0.00105493} & \num{89.0359} \\
93 & \num{0.000546502} & \num{0.00103276} & \num{89.0369} \\
94 & \num{0.000534936} & \num{0.0010113} & \num{89.038} \\
95 & \num{0.000523734} & \num{0.000990519} & \num{89.0389} \\
96 & \num{0.00051288} & \num{0.000970386} & \num{89.0399} \\
97 & \num{0.000502359} & \num{0.000950876} & \num{89.0409} \\
98 & \num{0.000492159} & \num{0.000931963} & \num{89.0418} \\
99 & \num{0.000482267} & \num{0.000913623} & \num{89.0427} \\
100 & \num{0.00047267} & \num{0.000895833} & \num{89.0436} \\
\end{longtable}

\section*{Caveats}
$\sigma_{CaF\text{-}He} = \num{1e-14}\,\text{cm}^2$ is a generic order-of-magnitude value for
molecule--He pairs used across the buffer-gas-beam literature; the true CaF--He elastic cross section
has energy-dependent structure (see e.g.\ Londo\~no \& P\'erez-R\'ios, \textit{J.\ Chem.\ Phys.}
\textbf{162}, 234304 (2025)) and could differ by a factor of a few. The near-field $z$-dependence
(item 6 in the formulas) is a standard free-jet scaling normalized to reproduce the paper's own $Re/2$
total, not a first-principles result for this specific cell. Background pressure is taken as the
paper's own quoted steady-state operating value ($\num{1e-7}$ mbar); the transient gas/plasma load from
the ablation pulse itself is not included.

\end{document}
